\documentclass{article}

\usepackage[swedish]{babel}
\usepackage{amsmath}
\usepackage{amssymb}

\title{Bevis av existens av Maclaurin utveckling}
\author{Marcus Persson}
\date{\today}

\begin{document}

\maketitle

\section*{\underline{Lemma 1: Cauchy's medelvärdessats}}

Antag att $g$ och $h$ är kontinuerliga på intervallet
$I=[a,b]$, där $a<b$, och deriverbara på $(a,b)$.
Då existerar minst ett $\xi \in (a,b)$ sådant att

\begin{equation*}
    (g(b)-g(a))h'(\xi)
    =
    (h(b)-h(a))g'(\xi).
\end{equation*}

\section*{\underbar{Bevis: existens av Maclaurin utveckling}}

Låt $g,h$ vara två kontinuerliga funktioner på intervallet $I=[0,x]$, där $0<x$, och deriverbara på $(0,x)$ där följande gäller:

\begin{equation}
    g(t) = \sum_{k=0}^{n}\frac{f^{(k)}(t)}{k!}(x-t)^k
\end{equation}

\begin{equation}
\begin{aligned}
g'(t)
&=
\sum_{k=0}^{n}\frac{f^{(k+1)}(t)}{k!}(x-t)^k
-
\sum_{k=1}^{n}\frac{f^{(k)}(t)}{(k-1)!}(x-t)^{k-1}\\
&=
\sum_{k=0}^{n}\frac{f^{(k+1)}(t)}{k!}(x-t)^k
-
\sum_{k=0}^{n-1}\frac{f^{(k+1)}(t)}{k!}(x-t)^k\\
&=
\frac{f^{(n+1)}(t)}{n!}(x-t)^n.
\end{aligned}
\end{equation}

\begin{equation}
    h(t) = (x-t)^{n+1}
\end{equation}

\begin{equation}
    h'(t) = -(n+1)(x-t)^n
\end{equation}

\begin{equation}
    g(0) = p_{n}(x)
\end{equation}

\begin{equation}
    g(x) = f(x)
\end{equation}

\noindent Vi använder oss av Cauchy's medelvärdessats där $\xi = \theta x$ och $0 < \theta < 1$, då följer:

\begin{equation}
\begin{gathered}
    (g(x)-g(0))h'(\xi)
    =
    (h(x)-h(0))g'(\xi)\\
    \Leftrightarrow\\
    (f(x)-p_n(x))(-(n+1)(x-\xi)^n)
    =
    -x^{n+1}\frac{f^{(n+1)}(\xi)}{n!}(x-\xi)^n\\
    \Leftrightarrow\\
    f(x)-p_n(x)
    =
    \frac{f^{(n+1)}(\xi)}{(n+1)!}(x-\xi)^{n+1}\\
    \Leftrightarrow\\
    f(x)
    =
    p_n(x) + \frac{f^{(n+1)}(\xi)}{(n+1)!}(x-\xi)^{n+1}
\end{gathered}
\end{equation}

Eftersom $\xi = \theta x$ och $0<\theta<1$ $\exists$ Maclaurin utveckling. 

\begin{flushright}
    $\blacksquare$
\end{flushright}

\end{document}